\section{The Gate as a System}
\label{sec:system}

\subsection{Accuracy--cost tradeoff}

On the frozen test split ($n=240$): small-only accuracy \textbf{0.588};
big-only (\texttt{gpt-5.5}) \textbf{0.917}; big relayed through the small
model's paraphrase 0.879 (relay tax 1--4 points). Every gate beats random
escalation at every operating point (Figure~\ref{fig:tradeoff}); the
balanced tier at a 33\% escalation budget reaches \textbf{0.779--0.787},
recovering $\sim$58\% of the small$\rightarrow$big gap at one third of the
big-only cost (\$1.12 per 100 queries big-only).

Figure~\ref{fig:tradeoff} also draws the baseline the type-shortcut audit
(\S\ref{sec:audit}) demands: a \emph{pool-oracle} router that escalates by
each pool's calibration-split base fail rate (true type labels, no
instance-level information, no test leakage). It captures $+0.042$ of the
gate's $+0.054$ area over random --- the system-level counterpart of the
AUC decomposition (0.715 of 0.821), and slightly harsher: in-distribution,
the instance-level residual the internal signal adds on top of type
routing is $+0.012$ area ($\sim$2 points at the balanced budget) --- not
significant at $n{=}240$ (paired bootstrap over test queries, $B{=}10^4$:
95\% CI $[-0.003, +0.027]$, two-sided $p{=}.13$); in distribution, the
gate is statistically on par with the pool oracle. The gate's case over
type routing therefore rests not on an in-distribution margin but on
transfer: the type shortcut is exactly what
collapses under LOPO (\S\ref{sec:audit}), while the mid-layer signal
survives it.

A \emph{trained} query-feature router closes this argument with data. We
fit a RouteLLM-style external router \citep{ong2024routellm} under
identical data conditions --- TF-IDF (word 1--2\,gram + character
3--5\,gram) $\to$ logistic regression on the calibration split's text
labels, seeing only the query surface. In-mix it lands \emph{below} even
the pool oracle (OOF AUC .669 vs.\ .715; tradeoff area $+.040$ vs.\ the
probe's $+.054$, mid-layer $+.064$, $\ptrue$-post $+.068$ on the same
test set and frozen expert answers); a 24-configuration sweep over
feature space and regularization tops out at .689, so the shipped
configuration is the query surface's ceiling, not a tuning artifact. Under LOPO it collapses entirely
(held-out-pool AUC .38--.57 vs.\ the mid-layer probe's .93 on math) and
its mean score on the 100\%-fail trap pool is .230 --- it \emph{misses}
the pool a router most needs to catch. The training receipt makes the
point starkly: at the argmax threshold the router's out-of-fold accuracy
on the calibration split, .678, \emph{exactly matches the majority-class
rate} (test .613 vs.\ majority .588) --- 360 queries of surface text buy
a weak ranking signal (train log-loss .38 vs.\ out-of-fold .59) and no
usable classifier. The internal probes pass the same accounting: at the
same threshold, same labels, and same logistic-regression machinery, the
text probe's out-of-fold accuracy is .772 (majority .678) and .779 on
test (majority .588), and the live audio mid-layer probe reaches
.764/.800 against .592/.512 --- real classifiers, not merely rankers;
the difference is purely the input representation
(Figure~\ref{fig:receipt}).

\begin{figure}[t]
  \centering
  \includegraphics[width=\linewidth]{receipt_compare}
  \caption{The same training receipt for the external router and the two
  internal probes (identical labels, identical logistic-regression
  machinery). (a)~Accuracy at the 0.5 threshold on the calibration
  (out-of-fold) and frozen test splits, against each split's
  majority-class baseline (dashed): the router never clears majority;
  both probes clear it by 9--29 points. (b)~Train vs.\ eval log-loss:
  the audio probe generalizes cleanly (.278 train / .446 test, the
  best-calibrated signal); the text probe memorizes in-sample
  (4096 dims $\gg$ 360 queries) yet its held-out accuracy holds --- all
  reported numbers are out-of-fold or test.}
  \label{fig:receipt}
\end{figure}

Nor is the baseline data-starved. The identical recipe trained on
RouterBench \citep{hu2024routerbench} --- 36{,}497 public prompts,
routing \texttt{mixtral-8x7b} $\to$ \texttt{gpt-4} (escalation base rate
.432) --- reaches out-of-fold AUC .710 in-domain (accuracy .660 vs.\
majority .568; deferral area over random $+.033$): a hundred times the
data buys .04 of AUC, landing at our pool-oracle's level and still short
of the probe's .828. Its leave-one-benchmark-out transfer sits at chance
on the format-disjoint benchmarks (HellaSwag .502, GSM8K .509,
Winogrande .498), reproducing our LOPO collapse on public data, and our
calibration-trained router transfers to RouterBench \emph{below} chance
(AUC .440). The transfer failure is symmetric: RouteLLM's
\emph{released} preference-trained routers (BERT and
matrix-factorization, $\sim$100k GPT-4-vs-mixtral preference pairs),
scored zero-shot on our pool by their own inference code, reach test AUC
only .523/.533 --- below even the 360-sample same-data router --- and
both rank the 100\%-fail trap pool \emph{below} the ordinary hard pools.
The audio modality narrows but does not close the gap: against the live
gate's audio labels the same recipe reaches OOF AUC .743 / test .814
(gold text; the talker's own ASR transcript costs it only
$-$.005--.009) --- stronger than on text labels, because the audio
channel makes hard pools fail more uniformly and the type shortcut buys
more --- yet the audio mid-layer probe still leads every readout
(.843/.879 AUC, accuracy .800 vs.\ .696), and the router still requires
the full utterance plus transcription, so it cannot fire mid-stream.
Query-feature routing is the type shortcut and nothing else; the
standardized LLMRouterBench protocol \citep{llmrouterbench2026} remains
future work, but the information-source question is settled under
identical data, the ceiling replicates at $100\times$ the data on a
public benchmark, and routing knowledge fails to port across talkers in
either direction.

\begin{figure}[t]
  \centering
  \includegraphics[width=.85\linewidth]{tradeoff}
  \caption{Accuracy vs.\ escalation rate, frozen test split: gated
  escalation vs.\ the random-escalation and pool-oracle (type-only)
  baselines between the small-only and big-only endpoints. The oracle
  captures most of the gate's in-distribution lift ($+0.042$ of $+0.054$
  area); the gate's advantage over type routing is transfer, not this
  margin (\S\ref{sec:audit}). See also Figure~\ref{fig:tradeoffptrue} for
  the $\ptrue$ variants.}
  \label{fig:tradeoff}
\end{figure}

\begin{figure}[t]
  \centering
  \includegraphics[width=.85\linewidth]{tradeoff_ptrue}
  \caption{Same tradeoff for the $\ptrue$ signals (pre-answer and
  post-draft).}
  \label{fig:tradeoffptrue}
\end{figure}

\subsection{Against the model's own thinking mode}
\label{sec:thinking}

MiniCPM-o~4.5 ships a built-in \texttt{enable\_thinking} reasoning mode ---
self-escalation in place. Table~\ref{tab:policy} compares deployment
policies with measured latencies. \textbf{Gated-cloud escalation dominates
all-THINK on both axes at every escalation rate.} The reason is mechanistic,
not incidental: the gate predominantly flags knowledge/trap failures, which
thinking cannot fix --- thinking helps math (fail .187 $\rightarrow$ .127)
but \emph{hurts} chat (.207 $\rightarrow$ .233) --- so external escalation
is necessary, not merely better.

\begin{table}[t]
\centering\small
\caption{Policy table, frozen test split ($n=240$). Latency from per-pool
measured medians; think and cloud latencies per-query.}
\label{tab:policy}
\begin{tabular}{lcccc}
\toprule
policy & acc & lat.\ mean (s) & P50 & P95 \\
\midrule
fast-only & .588 & 3.0 & 3.5 & 4.2 \\
all-THINK (o4.5's own) & .637 & 22.4 & 17.2 & 60.1 \\
gated-think @.33 & .613 & 12.1 & 3.6 & 47.6 \\
gated-cloud @.15 & .688 & 5.3 & 3.6 & 10.1 \\
\textbf{gated-cloud @.33} & \textbf{.787} & \textbf{6.5} & 3.6 & 20.8 \\
gated-cloud @.50 & .858 & 7.0 & 3.6 & 24.4 \\
\bottomrule
\end{tabular}
\end{table}

\subsection{The fork: deciding before the first token}
\label{sec:fork}

The mid-layer probe is one 4096-d dot product on a hidden state the forward
pass already computes: forking a decision branch at layer $k$ costs
microseconds and never stalls the main forward. The design question is at
which depth the signal is sufficient. Measured on H100 (CUDA-synchronized,
truncated-forward): the L22 decision lands in \textbf{20\,ms (text) /
45\,ms (audio)}, vs.\ time-to-first-token of 36 / 68\,ms --- the decision
exists \emph{before decoding starts} and sits far inside the 200--300\,ms
turn-taking budget \citep{stivers2009universals}. Per-layer curves
(Figure~\ref{fig:fork}) show L22's hidden state is available at 57\% of
prefill wall-time (61\% audio) while L22 is simultaneously the in-mix
quality peak (OOF .866 $>$ final layer .835): deciding early costs nothing,
and the cloud call can launch while the remaining $\sim$40\% of prefill
plus the whole decode still runs. Forking much earlier is not free:
pre-50\%-depth layers are both weaker (L10 $\approx$ .794) and
modality-specific (\S\ref{sec:xmodal}); the fork belongs at $\sim$55--60\%
depth. For comparison, $\ptrue$-pre costs 39 / 67\,ms (an extra short
prefill) and $\ptrue$-post costs a full draft (1.9 / 3.5\,s).

\begin{figure}[t]
  \centering
  \includegraphics[width=.85\linewidth]{fork_pareto}
  \caption{Fork-depth Pareto: per-layer truncated-forward wall-time vs.\
  gate quality. The deployed L22 sits at the quality peak while its hidden
  state is available at $\sim$57--61\% of prefill wall-time.}
  \label{fig:fork}
\end{figure}

\subsection{Does the cloud round-trip hide inside the talker's floor time?}
\label{sec:overlap}

If the gate fires at prefill and the cloud call launches immediately, the
expert result is ready before the local answer finishes speaking for
40--58\% of the whole test mix (audio input buys more time than text). But
\emph{gate-escalated} queries skew toward short local answers (trap and
knowledge pools; trap overlap $\approx$ 0), so pure overlap covers only
$\sim$20--31\% of escalations at a 33\% budget (Figure~\ref{fig:overlap}).
The residual wait after the local answer ends is P50 2--3\,s --- one stall
sentence --- but P90 $\sim$27\,s, dominated by the expert's own reasoning
latency. This argues for a fast-expert tier or streamed partial results in
step 2, and it quantifies the stall-phrase budget an injection protocol
must cover.

\begin{figure}[t]
  \centering
  \includegraphics[width=.85\linewidth]{overlap}
  \caption{Cloud round-trip vs.\ local answer duration: probability the
  expert result arrives before the local answer ends (+slack), whole mix
  and gate-escalated subset, text and audio.}
  \label{fig:overlap}
\end{figure}
